\documentclass{beamer}

\usepackage{amsmath,amssymb,amsfonts,dcolumn,color,graphicx,graphics,setspace,latexsym,setspace,lscape,subfigure,placeins,epsfig,hyperref}
\usepackage{eulervm}
\usepackage[utf8]{inputenc}

\usetheme{Kalgan}


\setbeamercovered{highly dynamic}
\setbeamersize{text margin left=10pt}

\newcommand{\be}{\begin{equation}}
\newcommand{\ee}{\end{equation}}
\newcommand{\lb}{\left}
\newcommand{\rb}{\right}



\begin{document}

\title{Linux WTF?}
\subtitle{Eine Einführung in das Betriebssystem Linux}
\author{Benedikt Schmitz}
%\institute[FUG]{FiwiAn UG}
\date{}



	\begin{frame}[plain]
	  \titlepage
	\end{frame}
	\begin{frame}
  		\frametitle{Outline}
		\tableofcontents
	\end{frame}
%%%%%%%%%%%%%%%%%%%%%%%%%%%%%
%%%% Inhalte 
%%%%%%%%%%%%%%%%%%%%%%%%%%%%%
	\section{Linux -- kann man das Essen?}
	    \subsection{Installation}
	        \begin{frame}
          		\frametitle{Linux -- kann man das Essen?}
            		\begin{center}
                		\begin{minipage}{0.44\textwidth}
                		    \begin{itemize}
                		        \item Installation
                		        \item Open Source
                		    \end{itemize}
                		\end{minipage}%
            		\end{center}
        	\end{frame}
        	
	        \begin{frame}
          		\frametitle{Installation via USB}
          		-- welche Einstellungen bei der Installation? Swap, zweite platte für /home
        		\begin{minipage}{0.44\textwidth}
        		
        		\end{minipage}%
        		\begin{minipage}{0.54\textwidth}
        		
        		\end{minipage}
        	\end{frame}
        	
	    \subsection{Open Source? -- Sinn!?}
	        \begin{frame}
          		\frametitle{Open Source/Free Software}
          		-- Vorteile/Nachteile?
        		\begin{minipage}{0.44\textwidth}
        		
        		\end{minipage}%
        		\begin{minipage}{0.54\textwidth}
        		
        		\end{minipage}
        	\end{frame}
        	
    \section{Linux $\leftrightarrow$ Windows: Unterschiede...}
        \begin{frame}
      		\frametitle{Linux $\leftrightarrow$ Windows: Unterschiede...}
        		\begin{center}
            		\begin{minipage}{0.44\textwidth}
            		    \begin{itemize}
            		        \item ... Konzeptionell
            		        \item ... Ordnerbaum
            		        \item ... Nutzer
            		        \item ... Installation von Software
            		    \end{itemize}
            		\end{minipage}%
        		\end{center}
    	\end{frame}
        \subsection{... Konzeptionell}
            \begin{frame}
          		\frametitle{Kernel/Hardware/Userland}
        		\begin{minipage}{0.44\textwidth}
        		
        		\end{minipage}%
        		\begin{minipage}{0.54\textwidth}
        		
        		\end{minipage}
        	\end{frame}
        	
        \subsection{... Ordnerbaum}
            \begin{frame}
          		\frametitle{Ordnerhierarchie/backups}
        		\begin{minipage}{0.44\textwidth}
        		
        		\end{minipage}%
        		\begin{minipage}{0.54\textwidth}
        		
        		\end{minipage}
        	\end{frame}
        	
        \subsection{... Nutzer}
            \begin{frame}
          		\frametitle{User/root}
        		\begin{minipage}{0.44\textwidth}
        		
        		\end{minipage}%
        		\begin{minipage}{0.54\textwidth}
        		
        		\end{minipage}
        	\end{frame}
        	
        \subsection{... Installation von Software}
            \begin{frame}
          		\frametitle{Installation / Repositories}
        		\begin{minipage}{0.44\textwidth}
        		
        		\end{minipage}%
        		\begin{minipage}{0.54\textwidth}
        		
        		\end{minipage}
        	\end{frame}

    \section{Wie verwendet man das dann?}
        \begin{frame}
  		\frametitle{Wie verwendet man das dann?}
    		\begin{center}
        		\begin{minipage}{0.44\textwidth}
        		    \begin{itemize}
        		        \item Beispiel: Thunderbird
        		        \item Beispiel: Installation
        		        \item Beispiel: eduroam?
        		        \item Beispiel: wine
        		    \end{itemize}
        		\end{minipage}%
    		\end{center}
    	\end{frame}
    	
        \subsection{Beispiel: Thunderbird}
            \begin{frame}
          		\frametitle{Einrichtung TB}
        		\begin{minipage}{0.44\textwidth}
        		
        		\end{minipage}%
        		\begin{minipage}{0.54\textwidth}
        		
        		\end{minipage}
        	\end{frame}
        	
        \subsection{Beispiel: Installation Chromium o.Ä.}
            \begin{frame}
          		\frametitle{Installation chromium oder anderer Software}
        		\begin{minipage}{0.44\textwidth}
        		
        		\end{minipage}%
        		\begin{minipage}{0.54\textwidth}
        		
        		\end{minipage}
        	\end{frame}
        
        \subsection{Beispiel: eduroam?}
            \begin{frame}
          		\frametitle{eduroam}
        		\begin{minipage}{0.44\textwidth}
        		
        		\end{minipage}%
        		\begin{minipage}{0.54\textwidth}
        		
        		\end{minipage}
        	\end{frame}
        	
        \subsection{Beispiel: wine}
            \begin{frame}
          		\frametitle{wine -- falls es nicht ohne Windows geht.}
        		\begin{minipage}{0.44\textwidth}
        		
        		\end{minipage}%
        		\begin{minipage}{0.54\textwidth}
        		
        		\end{minipage}
        	\end{frame}
        	
    \section{Hallo! Ich habe eine Problem...}
        \begin{frame}
      		\frametitle{Hallo! Ich habe eine Problem...}
    		\begin{center}
        		\begin{minipage}{0.44\textwidth}
        		    \begin{itemize}
        		        \item Das Ubuntuwiki
        		        \item Terminalbasics
        		    \end{itemize}
        		\end{minipage}%
    		\end{center}
    		\hfill \includegraphics[scale=0.2]{bilder/ohno.png}
    	\end{frame}
    	
        \subsection{Ubuntuwiki}
            \begin{frame}
          		\frametitle{Das Ubuntuwiki}
        		\begin{minipage}{0.44\textwidth}
        		
        		\end{minipage}%
        		\begin{minipage}{0.54\textwidth}
        		
        		\end{minipage}
        	\end{frame}
        	
	    \subsection{Terminalbasics}
    	    \begin{frame}
          		\frametitle{Das Terminal -- BASH}
          		-- kill, xkill 
                -- top 
                -- sudo apt-get ...
                -- scp? mc?    	
        		\begin{minipage}{0.44\textwidth}
        		    
        		\end{minipage}%
        		\begin{minipage}{0.54\textwidth}
        		
        		\end{minipage}
        	\end{frame}
        	
\end{document}
